% Last Updated: March 1, 2024
\documentclass[0-main]{subfiles}

\begin{document}
\ifSubfilesClassLoaded{
	% \title{\tbf{Untitled}}
	% \author{Cong Wen, xx/20xx}
	% \maketitle
	\tableofcontents
}{}
% ----------------------------------------------------------------------------------------------------

\ssc{Automorphic representations}

Ref:~\cite{Bump1997}, \cite{GH2024}

\sss{(Old treatment) On modular forms}

\spg{Sectional Goal}
\begin{itm}
	\item[\TODO] Do the following work
	\begin{itm}
		\item[\TODO] Hijikata--Pizer--Shemanske\cite{HPS1989Basis} resolved the basis problem for $S_{k}(\Gma_{1}(N),\psi)$ by decomposing them into twists of primitive neben space and theta spaces, therefore provided a way to concretely deal with them. Collect such classification results.
		\item[\TODO] Learn the treatment of modular forms in adelic settings, consider to use Zhang's thesis\cite{Zhang2023Derived}, or Gelbart\cite{Gelbart1975}.
	\end{itm}
\end{itm}


\sss{(obsolete understanding) Weil representations}

\spg{Sectional Goal}
\begin{itm}
	\item[\TODO] Find the answers for questions below:
	\begin{itm}
		\item[\TODO] Weil representations are powerful to deal with quadratic forms. Is there an analogue for cubic forms? One can use Manin's book\cite{MH1974}.
		\item[\TODO] (arithmetic) Siegel--Weil formulas have resolved deep local to global difficulties for quadratic equations. Then are the Hasse--Weil zeta functions for quadratic hypersurfaces completely known?
	\end{itm}
\end{itm}

\sss{The restriction problem}
Ref:~\cite{GGP2012}.

% ----------------------------------------------------------------------------------------------------
\ifSubfilesClassLoaded{
	\bibliographystyle{amsalpha}
	\bibliography{/Users/wenc/Documents/Math/universal-ref}
}{}
\end{document}
